\section{Introduction}
\label{sec:tpc40-introduction}

The last three layers of the TPC program separate three different
obstructions at an orbit--sliced prime--M\"obius energy gate.  TPC--37
closes the zero coordinate faces and the proper multiplicative faces, but
the full equality face retains the original physical vector with
coefficient \(1-o(1)\).  TPC--38 computes the target--synthesis spectrum
and shows that multiplicative completion produces a coherent joint--alias
bucket rather than the single equality column.  TPC--39 constructs the
missing additive quotient projector: exactly \(L+1\) twists recover the
zero alias, and no smaller arbitrary phase bank does so universally
\citep{WangTPC37,WangTPC38,WangTPC39}.

That algebraic recovery does not estimate the output.  Its minimal
positive square energy has the folded Parseval form
\begin{equation}
 \|\cZ_0\|^2+
 \sum_{s=1}^{L}\|\cZ_s+\cZ_{s-q}\|^2,
 \qquad q=L+1.
 \label{eq:tpc40-intro-folded-energy}
\end{equation}
In particular, \(\|\cZ_0\|^2\) is present as a positive summand.  At the
terminal layer, the alias phase is also independent of the physical row.
It therefore does not oscillate the fixed--alias cross--row quadratic
form, whose coefficients contain genuine two-- and four--M\"obius
correlations.

This suggests a second phase coordinate.  One could assign phases to the
physical rows, average the resulting energies, and hope to replace a
coherent row Gram by an almost diagonal one.  The difficulty is that the
same bank must still recover the original coherent row sum.  The present
paper determines the exact tradeoff between those two requirements.

\subsection{Main finite theorem}

Let \(R\) be the number of active rows, let \(\Phi\) be a unimodular
\(m\times R\) phase matrix, and let \(w\) be a positive probability
vector on its measurements.  If \(a\) recovers the coherent sum, define
\begin{equation}
 \cA(a,w)=\sum_{t=1}^{m}\frac{|a_t|^2}{w_t},
 \qquad
 G(\Phi,w)=\Phi^*\diag(w)\Phi.
 \label{eq:tpc40-intro-A-G}
\end{equation}
Our core result is the matrix inequality
\begin{equation}
 \boxed{
 G(\Phi,w)\succeq
 \frac1{\cA(a,w)}\one\one^*.}
 \label{eq:tpc40-intro-Loewner}
\end{equation}
It is stronger than an operator--norm estimate: every positive bank
energy that supports universal recovery already contains a fixed fraction
of the coherent target energy.  Since \(G\) has unit diagonal, one gets
the sharp alternatives
\begin{equation}
 \cA=1\Longrightarrow G=\one\one^*,
 \qquad
 G=I_R\Longrightarrow\cA\ge R.
 \label{eq:tpc40-intro-rigidity}
\end{equation}
Partial Fourier banks attain every integer point
\begin{equation}
 (\cA,\lambda_{\max}(G))=(r,R/r),
 \qquad 1\le r\le R,
 \label{eq:tpc40-intro-frontier}
\end{equation}
so the product law is not an artifact of the proof.

\subsection{Why this is the correct next layer}

The theorem does not claim that the physical M\"obius vector is a worst
case vector for \(G\).  It rules out a narrower and precisely declared
method class: coefficient--blind, equality--preserving phase banks that
are estimated through a positive weighted square energy.  Within that
class, row decorrelation is paid for by exactly the recovery
amplification, while the atomic diagonal is invariant.  A large sieve or
tight--frame statement about the phase geometry cannot alter this
conservation law.

This is useful even though it is an obstruction.  Before this paper,
adding row signs or a second Fourier coordinate remained a plausible
formal escape from the fixed--alias cross--row term in TPC--39.  The sharp
Loewner theorem closes that escape without making any conjecture about
M\"obius signs.  What remains is necessarily coefficient specific: one
must show that the actual arithmetic vector avoids the coherent mode, or
estimate a signed bilinear functional before replacing it by positive
energy.

\subsection{Guide to the paper}

\Cref{sec:tpc40-interface} records the inherited alias and endpoint
ledgers.  The weighted phase formalism, Loewner theorem, equality cases,
and sharp Fourier families occupy the next four sections.  We then tensor
the result with quotient tomography, treat coupled row--alias phases, and
specialize the conservation law to the physical endpoint.  A comparison
with multiplicative superwindows and with the nearest frame, large--sieve,
and M\"obius results keeps the scope explicit.  The final sections state
the coefficient--specific next gate and the exact computational
certificate.

\subsection{Claim boundary}

All phase--frame and matrix results in this paper are unconditional.  The
physical consequence is a no--go for the stated universal positive--energy
method class.  We do not prove that the actual cross--row energy is large,
that every phase method fails, or that signed coefficient--specific
dispersion is impossible.  In particular, no affine Chowla estimate,
parity--breaking estimate, Hardy--Littlewood formula, or twin--prime
theorem is asserted.
